\begin{proposition}[Refuted surface assertion, with explicit example]\hypertarget{tag-01F8}{}\label{cdp:st:C.3.2}
\StatementRecord{cdp:st:C.3.2}{proved}{01F8}{the definitions and constructions stated in the proof}


The assertion that every secondary Kodaira surface has torsion in
$H_1(-;\Z)$ is false. There is a secondary Kodaira surface $Q$ with
\[
H_1(Q;\Z)=\Z,\qquad H^2(Q;\Z)=0,
\]
whose canonical bundle has holomorphic order six.
\end{proposition}


\begin{proof}

\begin{enumerate}[label={$\langle1\rangle$\arabic*.},leftmargin=*]
\item Let $H$ be the real Heisenberg group in coordinates $(x,y,z)$ with law
\[
(x,y,z)(x',y',z')
=\left(x+x',y+y',z+z'+\frac{xy'-yx'}2\right).
\]
Let $\Gamma\subset H$ be generated by
$a=(1,0,0)$, $b=(0,1,0)$ and $c=(0,0,1)$. Thus
\[
\Gamma=\{(x,y,z):x,y\in\Z,\ z-xy/2\in\Z\},
\qquad [a,b]=c.
\]
It is a discrete cocompact subgroup. On $H\times\mathbb R_t$ write
\[
\tau=e^{\pi i/3},\quad
\eta=dz-\frac{x\,dy-y\,dx}{2},\quad
\alpha=dx+\tau\,dy,\quad
\beta=\eta-\frac{dx+dy}{2}+i\,dt.
\]
These two complex one-forms and their conjugates are linearly
independent over $\C$. Declaring $\alpha,\beta$ to have type $(1,0)$
defines a left-invariant almost complex structure. It is integrable:
$d\alpha=0$ and $d\beta=-dx\wedge dy$ has type $(1,1)$, so the
ideal generated by the $(1,0)$ forms has no $(0,2)$ obstruction.
The structure descends to
$P=\Gamma\backslash H\times(\mathbb R/\Z)$.
The nonzero form $\alpha\wedge\beta$ is $\bar\partial$-closed and
trivializes $K_P$. The projection induced by $x+\tau y$ is a
holomorphic elliptic bundle over $\C/(\Z+\tau\Z)$, with nonzero
central Chern class $[a,b]=c$. Its first Betti number is three, so
$P$ is a primary Kodaira surface.

\item Define the Lie-group automorphism
\[
\sigma(x,y,z)=(-y,x+y,z-y/2).
\]
It preserves $\Gamma$: the difference between the new central
coordinate and half the product of the new horizontal coordinates
is
\[
(z-xy/2)+xy+y(y-1)/2.
\]
This is integral on $\Gamma$. The inverse preserves the lattice too,
and direct substitution gives $\sigma^6=1$. Moreover
$\sigma^*\alpha=\tau\alpha$ and $\sigma^*\beta=\beta$.
Consequently
\[
g:P\longrightarrow P,\qquad
g([x,y,z],t)=([\sigma(x,y,z)],t+1/6)
\]
is a holomorphic automorphism of order six. Every nonidentity power
is fixed-point-free because of its translation of the $t$-circle.
The quotient $Q=P/\langle g\rangle$ is therefore a compact complex
surface and $P\to Q$ is finite \emph{\'etale} of degree six.

\item The contractible manifold $H\times\mathbb R$ is its universal cover.
The deck group is $\Gamma\rtimes_\sigma\Z$, where the infinite
generator is the lift $(\sigma,t\mapsto t+1/6)$; its sixth power is
the original integral translation in $t$. Since
$\Gamma^{\rm ab}=\Z^2$ and the induced matrix is
\[
A=\begin{pmatrix}0&-1\\1&1\end{pmatrix},\qquad
\det(A-I)=1,
\]
abelianization gives $H_1(Q;\Z)=\Z$.
The finite cover and $e(P)=0$ give $e(Q)=0$. Poincar\'e duality then
gives $b_2(Q)=0$, and the universal coefficient theorem gives
$H^2(Q;\Z)=0$ because $H_1(Q;\Z)$ has no torsion.
Finally $g^*(\alpha\wedge\beta)=\tau\alpha\wedge\beta$.
Every holomorphic section of $K_P^m$ is a constant multiple of
$(\alpha\wedge\beta)^m$, by compactness. Such a section descends
precisely when $6$ divides $m$. Thus $K_Q$ has holomorphic order six.
The surface has $b_1=1$ and a primary Kodaira surface as a finite
\'etale cover, so it is secondary Kodaira. The matrix and lattice
calculations have been checked exactly by the attached program.

\end{enumerate}
\end{proof}